\subsection{Peran Tokoh Kasman Singodimedjo}

Kasman Singodimedjo ditambahkan sebagai anggota PPKI oleh Sukarno bersama Wiranatakoesoemah, Ki Hajar Dewantara, Sajuti Melik, Iwa Koesoema Soemantri, dan Achmad Soebardjo, sehingga jumlah anggota PPKI bertambah dari 21 menjadi 27 orang.

\textbf{Peran Utama dan Perjuangan:}
\begin{itemize}
\item \textbf{Pemersatu Bangsa:} Menjadi pemersatu antara golongan Islam dan nasionalis pada rapat PPKI 18 Agustus 1945. Beliau dipercaya oleh Sukarno dan Mohammad Hatta untuk meyakinkan Ki Bagoes Hadikoesoemo agar menerima usulan penghapusan tujuh kata dalam Piagam Jakarta demi menjaga keutuhan bangsa.
\item \textbf{Kepemimpinan Militer:} Menjabat sebagai Ketua BKR Pusat yang dibentuk melalui Dekrit Presiden pada 23 Agustus 1945. Selanjutnya memimpin TKR, dan saat menjabat Ketua KNIP, beliau membubarkan TKR untuk membentuk TNI, sehingga beliau dikenal sebagai inisiator lahirnya TNI.
\item \textbf{Pemerintahan dan Hukum:} Menjadi Ketua KNIP (parlemen) pertama, Jaksa Agung kedua yang mengeluarkan Maklumat Jaksa Agung Nomor 3 Tanggal 15 Januari 1946 tentang pembenahan organisasi kejaksaan dan penegakan negara hukum, serta menjabat Menteri Muda Kehakiman pada Kabinet Amir Sjarifuddin II.
\item \textbf{Konstituante dan Gelar Pahlawan:} Terpilih sebagai anggota Dewan Konstituante dari Partai Masyumi pada Pemilu 1955. Ditetapkan sebagai Pahlawan Nasional pada 8 November 2018 berdasarkan Keppres Nomor 123/TK/Tahun 2018.
\end{itemize}

\subsection{Pembukaan Undang-Undang Dasar Negara Republik Indonesia Tahun 1945}

UUD NRI Tahun 1945 adalah hukum dasar tertinggi di Indonesia. Pembukaan UUD NRI Tahun 1945 disahkan oleh PPKI pada 18 Agustus 1945 dan diundangkan dalam Berita Republik Indonesia Tahun II Nomor 7 pada 15 Februari 1946. Kedudukannya sebagai hukum dasar tertulis berfungsi sebagai norma hukum yang mengikat seluruh penyelenggara negara dan warga negara, serta sebagai norma dasar yang menjadi sumber tertinggi bagi seluruh produk hukum.

\subsubsection{Isi dan Tujuan Pembukaan UUD NRI Tahun 1945}

Pembukaan UUD NRI Tahun 1945 terdiri atas empat alinea yang memuat makna mendalam:
\begin{itemize}
\item \textbf{Alinea Pertama:} Mengandung pengakuan hak kodrat dan prinsip universal bahwa kemerdekaan adalah hak segala bangsa, serta penjajahan harus dihapuskan karena tidak sesuai dengan peri kemanusiaan dan keadilan.
\item \textbf{Alinea Kedua:} Merupakan konsekuensi logis dari pernyataan kemerdekaan sebagai hasil perjuangan bangsa Indonesia untuk menentukan nasib sendiri guna mewujudkan cita-cita negara yang merdeka, bersatu, berdaulat, adil, dan makmur.
\item \textbf{Alinea Ketiga:} Memuat pengakuan nilai religius dan moral bahwa kemerdekaan serta pembentukan Negara Indonesia merupakan rahmat dari Tuhan Yang Maha Kuasa di samping hasil perjuangan bangsa.
\item \textbf{Alinea Keempat:} Memuat tujuan khusus negara (melindungi segenap bangsa dan seluruh tumpah darah Indonesia, memajukan kesejahteraan umum, mencerdaskan kehidupan bangsa) dan tujuan umum (ikut melaksanakan ketertiban dunia), ketentuan UUD negara, bentuk negara republik yang berkedaulatan rakyat, serta dasar filsafat negara (Pancasila).
\end{itemize}

\textbf{Makna dan Tujuan Pembukaan UUD 1945:}
\begin{itemize}
\item Menjadi sumber motivasi, aspirasi, cita-cita hukum, dan moral bangsa.
\item Memiliki nilai universal (dijunjung tinggi oleh bangsa-bangsa di dunia) dan nilai lestari (menampung dinamika masyarakat).
\item Menegaskan pernyataan kemerdekaan, menetapkan cita-cita bangsa, mengukuhkan Proklamasi, dan menetapkan dasar negara.
\end{itemize}

\subsubsection{Pembukaan UUD NRI Tahun 1945 sebagai Pokok Kaidah Negara yang Fundamental}

Sebagai \textit{staatsfundamentalnorm}, Pembukaan UUD NRI Tahun 1945 memenuhi dua unsur mutlak:
\begin{enumerate}
\item \textbf{Unsur Terjadinya:} Disusun oleh pembentuk negara (PPKI) yang merupakan representasi resmi seluruh bangsa Indonesia.
\item \textbf{Unsur Isinya:} Memuat dasar tujuan negara, ketentuan UUD negara, bentuk negara, dan dasar filsafat negara.
\end{enumerate}

\textbf{Alasan Pembukaan UUD 1945 Tidak Dapat Diubah:}
\begin{itemize}
\item Dibentuk oleh pembentuk negara (PPKI). Lembaga negara yang dibentuk setelahnya hanya merupakan alat perlengkapan negara dengan kedudukan lebih rendah.
\item Merupakan tertib hukum tertinggi dan memuat syarat-syarat mutlak bagi keberadaan Negara Republik Indonesia.
\item Merupakan pengejawantahan dari Proklamasi Kemerdekaan 17 Agustus 1945 yang hanya terjadi satu kali dalam sejarah.
\end{itemize}

\subsubsection{Hubungan Pembukaan UUD NRI Tahun 1945 dengan Proklamasi Kemerdekaan}

Pembukaan UUD NRI Tahun 1945 dan Proklamasi Kemerdekaan 17 Agustus 1945 merupakan satu kesatuan utuh yang tidak dapat dipisahkan.
\begin{itemize}
\item \textbf{Sifat Hubungan:} Pembukaan UUD memberikan penjelasan (menegakkan hak kodrat dan moral), penegasan (perjuangan membebaskan diri dari penjajahan yang diridai Tuhan), dan pertanggungjawaban (penyusunan UUD dan dasar negara) terhadap Proklamasi.
\item \textbf{Fungsi Proklamasi:} Proklamasi bukan tujuan akhir, melainkan prasyarat atau jembatan emas untuk mencapai tujuan nasional yang tercantum dalam Pembukaan UUD NRI Tahun 1945.
\end{itemize}

\subsubsection{Hubungan Pembukaan UUD NRI Tahun 1945 dengan Pasal-Pasalnya}

Hubungan Pembukaan UUD NRI Tahun 1945 dengan Batang Tubuh (Pasal-pasal) bersifat kausal organis:
\begin{itemize}
\item Pokok-pokok pikiran Pembukaan UUD meliputi suasana kebatinan dan mewujudkan cita-cita hukum (\textit{rechtsidee}) yang dijabarkan secara normatif ke dalam pasal-pasal UUD.
\item \textbf{Empat Pokok Pikiran Pembukaan UUD 1945:}
  \begin{enumerate}
  \item Negara melindungi segenap bangsa Indonesia dan seluruh tumpah darah Indonesia berdasar atas asas persatuan (Pokok Pikiran Persatuan).
  \item Negara hendak mewujudkan keadilan sosial bagi seluruh rakyat Indonesia (Pokok Pikiran Keadilan Sosial).
  \item Negara yang berkedaulatan rakyat berdasarkan atas kerakyatan dan permusyawaratan atau perwakilan (Pokok Pikiran Kedaulatan Rakyat).
  \item Negara berdasarkan atas Ketuhanan Yang Maha Esa menurut dasar kemanusiaan yang adil dan beradab (Pokok Pikiran Ketuhanan dan Kemanusiaan).
  \end{enumerate}
\end{itemize}

\begin{table}[h]
  \centering
  \caption{Penyelenggaraan Negara Berdasarkan Konstitusi}
  \begin{tabular}{|c|l|l|l|l|}
    \hline
    \textbf{No.} & \textbf{Penyelenggaraan Negara} & \textbf{UUD NRI Tahun 1945} & \textbf{Konstitusi RIS} & \textbf{UUDS 1950} \\ \hline
    1. & Bentuk Negara & Kesatuan & Serikat/Federasi & Kesatuan \\ \hline
    2. & Bentuk Pemerintahan & Republik & Republik & Republik \\ \hline
    3. & Sistem Pemerintahan & Presidensial & Parlementer & Parlementer \\ \hline
    4. & Kedudukan Presiden & Kepala Negara \& Pemerintahan & Kepala Negara & Kepala Negara \\ \hline
  \end{tabular}
\end{table}

\subsection{Dinamika Pelaksanaan UUD NRI Tahun 1945}

Pelaksanaan UUD NRI Tahun 1945 mengalami perkembangan dan perubahan sesuai dengan kondisi politik kenegaraan.

\textbf{Masa Berlakunya UUDS 1950 (1950--1959):}
\begin{itemize}
\item Menerapkan sistem pemerintahan parlementer yang menyebabkan ketidakstabilan politik dan kerap terjadi jatuh bangun kabinet sebanyak 7 kali pergantian Perdana Menteri (M. Natsir, Sukirman, Wilopo, Ali Sastroamidjojo I, Burhanuddin Harahap, Ali Sastroamidjojo II, dan Juanda).
\end{itemize}

\textbf{Dampak Negatif Instabilitas Politik Masa UUDS 1950:}
\begin{enumerate}
\item Program pembangunan jangka panjang pemerintah tidak dapat berjalan dengan lancar.
\item Terjadi ketidakharmonisan internal dalam tubuh Angkatan Bersenjata pasca-Peristiwa 17 Oktober 1952.
\item Ketegangan sosial politik meningkat akibat masa kampanye Pemilu 1953--1955 yang berkepanjangan.
\item Timbulnya gerakan separatisme dan pemberontakan di daerah seperti PRRI di Sumatra dan Permesta di Sulawesi.
\item Kebuntuan sidang Badan Konstituante dalam menetapkan dasar negara baru.
\end{enumerate}

Kondisi krisis tersebut mendorong Presiden Sukarno untuk mengajukan konsep Demokrasi Terpimpin dan mengembalikan berlakunya UUD NRI Tahun 1945 melalui Dekrit Presiden.
